\section*{第 \ref{ch:g11:genes-and-disease} 章——遗传变异与疾病}

\begin{solution}{exo:g11:genes-and-disease:1}
\omterm{def:g11:genes-and-disease:genetic}{隐性}：只有两个\omterm{def:g10:universal-dna:gene}{等位基因}拷贝都存在时才出现疾病；
\omterm{def:g11:genes-and-disease:genetic}{显性}：一个拷贝即可。携带者带有一个\omterm{def:g11:genes-and-disease:genetic}{隐性}\omterm{def:g11:mutations:mutation}{突变}\omterm{def:g10:universal-dna:gene}{等位基因}，
但表现正常。
\end{solution}

\begin{solution}{exo:g11:genes-and-disease:2}
患病 $1/4$，携带者 $1/2$，两者皆非 $1/4$。
\end{solution}

\begin{solution}{exo:g11:genes-and-disease:3}
I-1 为 $Aa$，I-2 为 $Aa$（患病子女的父母），II-2 为 $aa$，III-2
为 $aa$。
\end{solution}

\begin{solution}{exo:g11:genes-and-disease:4}
男性只有一条 X \omterm{def:g11:cell-cycle-mitosis:chromosome}{染色体}：其上的\omterm{def:g11:genes-and-disease:genetic}{隐性}\omterm{def:g10:universal-dna:gene}{等位基因}会表达，
因为没有第二个拷贝可以掩盖它。女性则须两条 X \omterm{def:g11:cell-cycle-mitosis:chromosome}{染色体}上
都带有该\omterm{def:g10:universal-dna:gene}{等位基因}。
\end{solution}

\begin{solution}{exo:g11:genes-and-disease:5}
其发病风险取决于多个效应较小的\omterm{def:g10:universal-dna:gene}{等位基因}及环境的疾病：
2 型糖尿病、心脏病（还有哮喘、大多数
癌症）。
\end{solution}

\begin{solution}{exo:g11:genes-and-disease:6}
$aa \times AA$：子代均为 $Aa$，即健康的携带者。携带者子女
与携带者婚配：$1/4$ 患病，$1/2$ 为携带者，$1/4$ 为非携带者。
\end{solution}

\begin{solution}{exo:g11:genes-and-disease:7}
II-1 是两名携带者的健康子女，其为携带者的概率为
$2/3$。与人群中一名男性所生子女的风险：$2/3 \times 1/25 \times 1/4
= 1/150$.
\end{solution}

\begin{solution}{exo:g11:genes-and-disease:8}
每个子女的概率为 $1/2$。四个子女均不患病：$(1/2)^4 = 1/16$。
\end{solution}

\begin{solution}{exo:g11:genes-and-disease:9}
$X^hY \times X^HX^H$：女儿全部为 $X^HX^h$，均为健康的携带者；
儿子全部为 $X^HY$，均健康。携带者女儿与健康男性婚配：
其一半儿子患血友病，一半女儿为携带者。
\end{solution}

\begin{solution}{exo:g11:genes-and-disease:10}
一个\omterm{def:g10:universal-dna:gene}{等位基因}产生 400 和 200 两条条带；另一个\omterm{def:g10:universal-dna:gene}{等位基因}再加上缺失，
产生 300 和 200 两条条带。条带位于 400、300 和 200 处——无 600 处。
\end{solution}

\begin{solution}{exo:g11:genes-and-disease:11}
\omterm{def:g10:universal-dna:gene}{等位基因}在种群中的频率取决于其历史，
而不只取决于\omterm{def:g11:mutations:mutation}{突变}速率：欧洲的该\omterm{def:g10:universal-dna:gene}{等位基因}很早就已出现并扩散，
或许因为携带者具有某种优势（有人怀疑其对某种腹泻病有抗性）；
在东亚则未扩散。
\end{solution}

\begin{solution}{exo:g11:genes-and-disease:12}
每个子女都是独立事件：第二个子女患病的概率为 $1/4$。
接下来两个子女中至少一个患病：$1 - (3/4)^2 = 7/16$。
每次受孕时\omterm{def:g10:universal-dna:gene}{等位基因}都重新分配；前一个患病子女并不会降低
下一个子女的风险。
\end{solution}

\begin{solution}{exo:g11:genes-and-disease:13}
检测前，某女性为携带者的概率为 $1/25$；检测会漏掉
10\% 的携带者，因此在 1000 名女性中 40 名为携带者，其中 4 人检测为
阴性，而 960 名非携带者检测为阴性：携带者概率为
$4/964 \approx 1/240$。已知为携带者时的风险：$1/240 \times 1/4
\approx 1/960$。
\end{solution}

\begin{solution}{exo:g11:genes-and-disease:14}
由多个风险\omterm{def:g10:universal-dna:gene}{等位基因}降为少个、仍久坐：由 34\% 降至 9\%；由久坐改为积极运动、
仍有多个风险\omterm{def:g10:universal-dna:gene}{等位基因}：由 34\% 降至 12\%。可以改变的因素所带来的降幅，
几乎与无法改变的因素相同。风险评分有助于识别谁能从改变生活方式中获益最多，
而非作为
定论。
\end{solution}

\begin{solution}{exo:g11:genes-and-disease:15}
检测带来确定性——若结果为阴性，即摆脱五成的阴影；
若结果为阳性，则能规划人生与家庭——但
也有代价：阳性结果会使人提前数十年得知一种无法治愈的疾病，
并对情绪、保险以及同样携带该\omterm{def:g10:universal-dna:gene}{基因}的亲属产生影响。
检测结果并不会带来任何治疗，因此了解结果的唯一理由
在于本人的意愿；这也正是检测与否及何时检测
应由其本人决定的原因。
\end{solution}

\begin{solution}{pb:g11:genes-and-disease:1}
\textbf{1.} \omterm{def:g11:genes-and-disease:genetic}{隐性}：佐伊患病，而其父母均健康，
因此二人各携带一个\omterm{def:g11:genes-and-disease:genetic}{隐性}\omterm{def:g11:mutations:mutation}{突变}\omterm{def:g10:universal-dna:gene}{等位基因}。

\textbf{2.} 莱雅为 $Aa$，汤姆为 $Aa$，佐伊为 $aa$。

\textbf{3.} $AA$ 的概率为 $1/3$，$Aa$ 的概率为 $2/3$（他
健康，故排除 $aa$）。

\textbf{4.} 汤姆是携带者，因此其父母至少有一人是携带者；由于该代无患病者，
安娜从该父母处获得该\omterm{def:g10:universal-dna:gene}{等位基因}的概率为
$1/2$：$1/2$（设另一位父母为
非携带者）。

\textbf{5.} $1/25$，即人群中的频率。

\textbf{6.} $1/2 \times 1/25 \times 1/4 = 1/200$.

\textbf{7.} 患病者 $1/4$；携带者 $1/2$。

\textbf{8.} $2/3 \times 1/25 \times 1/4 = 1/150$.

\textbf{9.} 佐伊为 $aa$；其配偶为 $Aa$ 的概率为 $1/25$。
患病者：$1/25 \times 1/2 = 1/50$。携带者：每个孩子都从佐伊处获得 $a$，
因此为携带者的概率为 $24/25 \times 1 + 1/25 \times
1/2 \approx 0.98$。

\textbf{10.} 卢为携带者的概率为 $2/3$；其表亲
从安娜处获得 $a$（安娜为 $1/2$ 携带者）的概率为 $1/2$，即其
为携带者的概率为 $1/4$（忽略卡里姆）：$2/3 \times 1/4
\times 1/4 = 1/24$，为马克斯的四倍，因为该表亲
与卢同属一个家族，比陌生人更可能
携带同一\omterm{def:g10:universal-dna:gene}{等位基因}。

\textbf{11.} 佐伊：仅 600 和 300。汤姆：900、600 和 300。非携带者：
仅 900。

\textbf{12.} $AA$：他不是携带者；其子女的风险由
$1/150$ 降至零（排除新\omterm{def:g11:mutations:mutation}{突变}）。

\textbf{13.} 安娜为 $Aa$：风险为 $1 \times 1/25 \times 1/4 = 1/100$。

\textbf{14.} 在与卡里姆情况相似的 1000 人中，40 人为携带者，其中 4
人检测为阴性，960 名非携带者同样检测为阴性：$4/964 \approx
1/240$。风险：$1 \times 1/240 \times 1/4 \approx 1/960$。

\textbf{15.} $Aa$：健康的携带者，与其母亲相同。

\textbf{16.} 两人均为携带者的夫妇占 $(1/25)^2 = 1/625$；其
子女中四分之一患病：$1/2500$。

\textbf{17.} 恰好等于观察到的频率：携带者频率
与\omterm{def:g11:genes-and-disease:genetic}{隐性}遗传规律共同解释了这一频率。

\textbf{18.} 携带者频率为 $1/25$，患病者
频率为 $1/2500$，健康携带者所携带的\omterm{def:g11:mutations:mutation}{突变}\omterm{def:g10:universal-dna:gene}{等位基因}
数量比患病者携带的多，两者之比为 $2 \times 1/25$ 比 $2 \times
1/2500$，即一百比一。患病者死亡每代仅消除 1\% 的
拷贝；其余则由
携带者无形传递。

\textbf{19.} 该\omterm{def:g10:universal-dna:gene}{等位基因}的频率将会上升，但上升非常缓慢：
过去每代丢失的 1\% 拷贝如今被保留，因此
每代的变化量约为频率的 1\% 量级。

\textbf{20.} 两名携带者生育患病子女的风险为四分之一；患病子女的
健康同胞为携带者的概率为三分之二；凝胶
将概率转化为\omterm{def:g11:enzymes-and-phenotype:phenotype}{基因型}——排除了马克斯，证实了安娜，
令人安心但未能排除卡里姆，并确定胎儿为
健康的携带者。
\end{solution}
